\documentclass[11pt]{article}
\usepackage{amsmath,amssymb,amsthm,hyperref}
\newtheorem{theorem}{Theorem}
\newtheorem{lemma}[theorem]{Lemma}
\newtheorem{corollary}[theorem]{Corollary}
\title{A single middle vertex is isolated}
\author{human\_compression project}
\date{}
\begin{document}\maketitle

\paragraph{Scope and attribution.}
The original target is the quintic, girth-four, diameter-eight cage census
in Cambie, Goedgebeur, Jooken and Van den Eede,
\emph{On the order-diameter ratio of girth-diameter cages},
DMTCS 28:4, SOFSEM 2026 \#1, arXiv:2511.21144v3,
19 August 2026, Appendix B, Table 6, physical page 23.
The final parameter eight is a graph distance, not an edge-cycle
multiplicity. This proof uses only elementary graph arguments under the
hypotheses below. It imports no claimed numerical census, group action,
graph file or computational result. The independent-neighborhood and side
definitions are borrowed from the retained F009 defect-reduction work.
The present deduction and its application are a project assembly, with
no novelty or priority claim.

This note proves a conditional statement about complete graphs.
Its application to a partial family requires the certificates stated below.
The proof uses no computational census or graph file.

\paragraph{Ambient hypotheses.}
Let $G$ be a connected simple triangle-free $5$-regular graph on exactly
$32$ vertices. Suppose that
\[
 p_0,p_1,\ldots,p_8
\]
is a geodesic of length eight. Thus $d_G(p_0,p_8)=8$ and all of its
subpaths are shortest paths. We do \emph{not} assume that the global
diameter of $G$ is eight. Put
\[
 A=N(p_0),\quad B=N(p_4),\quad C=N(p_8),\quad
 I=A\cup B\cup C,\quad J=V(G)\setminus I,\quad F=G[J],
\]
and, on $K=J\setminus\{p_0,p_4,p_8\}$, define
\[
 L=\{x\in K:N(x)\cap A\ne\varnothing\},\quad
 R=\{x\in K:N(x)\cap C\ne\varnothing\},\quad
 M=K\setminus(L\cup R).
\]
These are definitions in the complete graph. The theorem concerns
\emph{exactly} $M=\{h\}$, not merely one named vertex missing side
neighbors in a partial graph.

\begin{lemma}[The geodesic side setup]\label{setup}
The set $I$ is independent of size $15$, $|J|=17$, $|K|=14$,
and $F$ has exactly five edges. The three centers are isolated in $F$.
The sets $L,R$ are disjoint, have size at least five, and have no edge
between them. No vertex of $B$ meets both $L$ and $R$.
For $D=G[R]$ and $x\in R$,
\[
 d_D(x)\leq |R|-4,                                  \tag{1}
\]
and the analogous bound holds on $L$.
\end{lemma}
\begin{proof}
The center distances are $4,4,8$. Neighborhood intersections or
edges between two of their neighborhoods would produce a center path
of length two or three. Each individual neighborhood is independent
by triangle-freeness. Hence $I$ is an independent $15$-set, and the
centers lie in $J$ and have all their neighbors in $I$.
The graph has $80$ edges, while the independent set $I$ has $75$
incident edges, all leading to $J$. Exactly five edges remain in $F$.

A vertex meeting both $A$ and $C$ would shorten the endpoint distance
to four; an $L$--$R$ edge would shorten it to five; a $B$-vertex
meeting both sides would shorten it to six. None is possible.
Every $a\in A$ needs four distinct neighbors in $L$ besides $p_0$,
so $|L|\geq4$. Equality closes a proper $K_{5,5}$ component between
$A$ and $\{p_0\}\cup L$, contradicting connectedness and order $32$.
Thus $|L|\geq5$, and similarly $|R|\geq5$.
For $x\in R$, choose a $C$-vertex adjacent to $x$. Its other three
side neighbors cannot neighbor $x$, by triangle-freeness. This gives
(1). Every argument used only the fixed geodesic, not global diameter.
\end{proof}

\begin{theorem}[Singleton-middle isolation]\label{main}
Under the ambient hypotheses, if $M=\{h\}$ then
\[
 d_F(h)=0,\qquad N_G(h)=B=N_G(p_4).
\]
In particular, the stratum with one middle vertex and a positive
internal middle degree is empty. The vertices $h,p_4$ are nonadjacent
false twins. No global-diameter-eight assumption is required.
\end{theorem}
\begin{proof}
Suppose, for a contradiction, that $d=d_F(h)>0$.
Because $h\in M$ meets neither $A$ nor $C$, and the centers are
isolated in $F$, all its $F$-neighbors lie in $L\cup R$.
If it had neighbors $x\in L$ and $y\in R$, the walk
\[
 p_0,A,x,h,y,C,p_8
\]
would give an endpoint path of length at most six, a contradiction.
After interchanging the endpoints and sides, all $d$ internal
neighbors of $h$ lie in $L$. Its other $5-d$ neighbors lie in $B$.
We keep the full range $1\leq d\leq5$; connectedness alone has
not been used to assume $d\leq4$.

Each $b\in B$ meets $p_4$ and has either three or four remaining
side neighbors, according as it does or does not meet $h$.
It meets just one side by Lemma~\ref{setup}. Write $B_L,B_R$
for these two parts of $B$. If $b\in B_R$ also met $h$, choose
an internal neighbor $x\in L$ of $h$ and a side neighbor $y\in R$
of $b$. Then
\[
 p_0,A,x,h,b,y,C,p_8
\]
would shorten the endpoint distance to at most seven. Hence
\[
 N_B(h)\subseteq B_L,\qquad N(h)\cap R=\varnothing.  \tag{2}
\]

Both parts $B_L,B_R$ are nonempty. If $B_R$ were empty, then
$R\cup C\cup\{p_8\}$ would have no edge to its complement:
there are no $L$--$R$ edges or $h$--$R$ edges. If $B_L$ were
empty, (2) would also give $N_B(h)=\varnothing$, and
$L\cup A\cup\{p_0,h\}$ would similarly be a proper component.
Both contradict connectedness. Put $b=|B_R|$. The $5-d$ distinct
$B$-neighbors of $h$ all lie in $B_L$, so
\[
 1\leq b\leq4,\qquad b\leq d.                     \tag{3}
\]

Put $n=|R|$ and $e=e(G[R])$. The side sizes sum to thirteen,
so $5\leq n\leq8$. Every $C$-vertex has four neighbors in $R$,
and, by (2), every $B_R$-vertex also has four neighbors there.
No other external neighborhood is available to a vertex of $R$.
The degree sum on $R$ gives
\[
 5n=20+4b+2e,
 \qquad e\leq5-d,\qquad\Delta(G[R])\leq n-4.       \tag{4}
\]
The edge inequality uses the five total $F$-edges, of which $d$
already touch $h$. Each of the $5+b$ rows supplied by $C\cup B_R$
is an independent four-subset of $G[R]$.

Equation (4) makes $n$ even, hence $n=6$ or $n=8$.
If $n=8$, then $e=10-2b$. Equations (3)--(4) imply
\[
 10-2b\leq5-d\leq5-b,
\]
so $b\geq5$, contradicting (3). This also covers the initially
allowed possibility $d=5$, without assuming it away.

It remains that $n=6$. Now
\[
 e=5-2b\geq0,
\]
so $(b,e)$ is $(2,1)$ or $(1,3)$.

If $(b,e)=(2,1)$, the side graph is one edge plus four isolates.
Each endpoint of that edge has internal degree one and consequently
needs four neighbors among $C\cup B_R$. The two endpoints together
need eight incidences with these seven rows. But each row is an
independent set and can contain at most one endpoint. Contradiction.

If $(b,e)=(1,3)$, the side graph is triangle-free of maximum degree
two, on six vertices and with three edges. It has no cycle, since
a cycle of length three is forbidden and a longer cycle uses more
than three edges. Its three path components have sizes
$(4,1,1)$, $(3,2,1)$ or $(2,2,2)$.
For $3P_2$ the independence number is three, so there is no row of
size four. For $P_4\sqcup2K_1$, every independent four-subset
contains both isolates. For $P_3\sqcup P_2\sqcup K_1$, every
independent four-subset contains its isolated vertex (and both
endpoints of the $P_3$). In either of the latter two cases an
isolate must neighbor all six distinct rows of $C\cup B_R$, giving
degree at least six. All alternatives are impossible.

This contradiction proves $d_F(h)=0$. All five neighbors of $h$
then lie in $I$, and, since $h\in M$, all lie in the five-set $B$.
Thus $N(h)=B=N(p_4)$, proving the theorem.
\end{proof}

\paragraph{Why the weaker hypothesis matters.}
A restored F009 family keeps its marked-root distance at least eight.
If its required graph contains a nine-distinct-vertex path between
those marked roots, that path is geodesic in every completion.
The theorem above therefore applies without first proving an all-pairs
diameter bound. Connectedness, degree five, order $32$ and
triangle-freeness remain essential hypotheses. In an actual selected
family they are separately enforced by the retained target contract.

\begin{corollary}[A sufficient partial-family emptiness certificate]
Let a partial family have the authentic full restored-family contract
and target $(32,5,4,8)$. Suppose required edges supply a nine-vertex
marked-root path and saturate its centers $p_0,p_4,p_8$ at degree five.
Let $A,B,C,I,J,K$ be their fixed sets. Suppose a supplied $h\in K$
satisfies the following \emph{universal} conditions:
\begin{enumerate}
\item no completion can give $h$ a neighbor in $A\cup C$;
\item every $x\in K\setminus\{h\}$ has a required neighbor in
$A$ or $C$;
\item some required edge joins $h$ to a vertex of $J$.
\end{enumerate}
Then the family has no selected completion.
\end{corollary}
\begin{proof}
In every selected completion the required path and root-distance
contract give the geodesic. Saturation fixes the center neighborhoods.
The first two conditions make $M$ exactly $\{h\}$, whereas the
last forces $d_F(h)>0$. Theorem~\ref{main} forbids this.
\end{proof}

One direct sufficient certificate for the first condition is that
every vertex of $A\cup C$ is already saturated in the required graph
and no required edge joins it to $h$. More economical certificates
could use separately proved triangle or marked-distance exclusions.
A presently absent edge alone is not such a proof. An unsuccessful
profile test is unsupported, not an emptiness theorem.

This gives a genuine selected-family \texttt{EMPTY} conclusion when
all hypotheses are certified, unlike the different M3 prior-union
coverage rule. It requires a distinct reviewed theorem/queue interface
before any adoption. It does not fabricate a native return, leaf
occurrence count, successful Source run, global class count, or
resolution of a reserved dispatch. No current family was inspected.

\paragraph{Relationship to the frozen M2 and M3 notes.}
The setup proof shows explicitly that its restrictions need only the
fixed length-eight geodesic. The frozen M2/M3 proofs also have their
center/side setup stated using global diameter eight; their template
existence proofs separately establish that diameter. This suggests a
further hypothesis weakening, but this note does \emph{not} silently
change their immutable statements, ordinary-class scopes or reviews.
Any strengthened M2/M3 edition must receive a separate exact proof
audit showing that no intermediate step used the global diameter.
The present singleton result needs no such strengthening or imported
M2/M3 classification theorem.
\end{document}
